\documentclass[11pt]{article}
\usepackage{amsmath,amssymb,amsthm,mathrsfs}
\usepackage[margin=1in]{geometry}
\usepackage[hidelinks]{hyperref}
\newcommand{\DD}{d'd''}
\newcommand{\MA}{\operatorname{MA}}
\newcommand{\C}{\mathcal C}
\newcommand{\G}{\mathcal G}
\newcommand{\GN}{\mathcal G_{\mathrm{net}}}
\newcommand{\RR}{\mathbb R}
\newcommand{\HH}{\mathscr H}
\newtheorem{theorem}{Theorem}
\newtheorem{lemma}[theorem]{Lemma}
\newtheorem{corollary}[theorem]{Corollary}
\newtheorem{proposition}[theorem]{Proposition}
\theoremstyle{remark}\newtheorem{remark}[theorem]{Remark}
\title{Uniform Gauss bounds and bounded-test tail control\\on the actual nonalgebraic analytic surface}
\author{Research working paper 407; unreviewed}
\date{9 October 2026}
\begin{document}
\maketitle
\begin{abstract}
The normalized continuous positive weights on the actual analytic
torus quotient have reference Gauss values between $-1/2$ and zero,
independently of their energy. We prove this by comparison on two
actual coordinate annuli, not by treating a varying Shilov section
as an analytic map. Fixed bounded-test gradient measures are
dominated by finite sums of actual bounded mixed Monge--Ampere
measures. The bounded Green calculus then yields a uniform
$T^{-1/2}$ bound for negative-part truncation errors on every
normalized energy sublevel, against each fixed bounded test's
Monge--Ampere measure. The bound passes to actual decreasing nets,
including poles. Thus these families are uniformly integrable
against every such fixed measure. No arbitrary-limit
representability, full energy-sublevel compactness, or original
equation-existence theorem is claimed. Original P16.7 remains open;
all new claims are unreviewed.
\end{abstract}

\section{Actual surface, class, and raw normalization}
Use
\[
 K=\mathbb C((s)),\quad |s|=e^{-1/8},\quad t=s^8,
 \quad X=(\mathbb G_m^2)^{\mathrm{an}}_K/
              \langle(t,2),(3,t)\rangle.
\]
This is the actual compact, separated, good, smooth, boundaryless
analytic quotient reconstructed earlier. Quotient geometry and
nonalgebraization are disclosed prior foundations, not reproved
here. Put $r_i=-\log|z_i|$, $\rho=|r|^2/2$ on a chart and
$\theta=\DD\rho$. Period changes of $\rho$ are affine unit
logarithms, so $\theta$ is global. The raw mass is $M=2$.
The all-stage continuous CLD class $\C$, its positive products
$B(a,b)=\theta_a\wedge\theta_b$, its bounded Green calculus and
concave energy $E$ are supplied in the full appendix. We retain
every continuous stage, rather than presume a first-stage identity.
Write $Q_T(f,g)$ for its gradient bilinear form against a positive
closed $(1,1)$ current $T$. Continuous differences of local psh
functions belong to this actual strong-gradient domain.

Let $\sigma(r)$ be the joint coordinate Gauss point, mapped to
$X$. It is one-periodic in each real coordinate: multiplication
by $(t,2)$ and $(3,t)$ shifts the coordinate norms by $(1,0)$ and
$(0,1)$, and the norm-one constants preserve Gauss seminorms.
Its image is the reference skeleton $S\simeq\RR^2/\mathbb Z^2$.
Use $\GN$ for actual decreasing continuous presentations with a
common finite lower energy bound, as in the closure proof406.
The net arguments needed below are proved here explicitly; that
paper is retained provenance, not an unexplained correctness grant.

\section{Two annulus comparisons give uniform reference bounds}
\begin{theorem}\label{gauss}
For $a\in\C$, put $g(r)=a(\sigma(r))$. Then $g$ is separately
$1/2$-Lipschitz in the real coordinates and
\[
 |g(r)-g(r')|\le\frac12\sum_{i=1}^2
      d_{\RR/\mathbb Z}(r_i,r_i').
\]
Every point $x$ of tropical coordinate $r$ satisfies $a(x)\le g(r)$.
Consequently $\sup_Xa=\max_Sa$; if $\sup_Xa=0$, then
\begin{equation}\label{refbound}
 -\tfrac12\le a|_S\le0.
\end{equation}
\end{theorem}
\begin{proof}
Fix $r_2$. Let $F=\mathscr H(\eta_{e^{-r_2}})$ be the completed
residue field of the actual $z_2$ Gauss point. The fibre is the
actual $F$-analytic curve $(\mathbb G_m)^{\mathrm{an}}_F$ with
coordinate $z_1$. More explicitly, its morphism to the torus
after extension to $F$ sets $z_2$ equal to its nonzero canonical
element in $F$; composing with the quotient and scalar projection
gives the actual pullback of $a$. The local weights $\rho+a$ are
classically psh: smooth CLD recognition and decreasing closure
recognize all continuous stages as in the appendix. Scalar
extension and analytic pullback preserve this property by GRv2
Proposition4.2(5)--(6).

On the whole curve the resulting weight has radial Gauss values
\[
 \Phi(q)=\tfrac12(q^2+r_2^2)+g(q,r_2).
\]
Local reference changes are pluriharmonic affine unit logarithms,
so these expressions glue to the actual classical subharmonic
weight on the coordinate curve. On an annulus with real radial
interval $[b,c]$, its two Shilov points are the endpoint Gauss
points. The affine interpolation of $\Phi(b)$ and $\Phi(c)$ is
pluriharmonic, being affine in $-\log|z_1|$. Apply the affinoid
maximum principle to the difference. It gives
\[
 \Phi(q)\le\frac{c-q}{c-b}\Phi(b)+
                 \frac{q-b}{c-b}\Phi(c),\quad b\le q\le c.
\]
Thus $q\mapsto q^2/2+g(q,r_2)$ is convex on $\RR$.
This uses actual annuli over $F$, not an analytic map defined by
varying a higher-dimensional Shilov section. Reversing the roles
of the two coordinates gives the same conclusion in the other
variable, on the other actual coordinate curve.

If $h$ is one-periodic and $q^2/2+h(q)$ is finite convex, every
subgradient $p$ of that convex function at $q$ satisfies
\[
 q-\tfrac12\le p\le q+\tfrac12.
\]
Indeed these are the secant slopes to $q-1$ and $q+1$, using
periodicity. Finite convex functions are locally absolutely
continuous, so $|h'|\le1/2$ almost everywhere and $h$ is
$1/2$-Lipschitz. Periodicity replaces real distance by circle
distance. Applying this one coordinate at a time proves the
displayed inequality and joint continuity of $g$. The maximal
circle distance is $1/2$, so $\operatorname{osc}g\le1/2$.

Lift $x$ to the torus and surround its tropical coordinate $r$
by a sufficiently small closed rectangular polyannulus inside
one injective quotient chart. Its Shilov boundary consists of the
four joint Gauss vertex points. The maximum principle for
$\rho+a$ bounds its value at $x$ by the largest vertex value.
Shrink the rectangle to $r$; continuity of $\rho+g$ gives
$a(x)\le g(r)$. The supremum assertion and\eqref{refbound} follow.
\end{proof}

\begin{proposition}\label{refmeasure}
The actual Radon measure $\theta^2$ is supported on $S$ and is
the pushforward of twice Lebesgue measure on $[0,1]^2$.
\end{proposition}
\begin{proof}
First work on the torus. A point $x$ outside its coordinate
Gauss skeleton has a Laurent polynomial $P=\sum_Ic_Iz^I$ with
\[
 |P(x)|<\max_I|c_Iz^I(x)|.
\]
Otherwise its seminorm on every Laurent polynomial is the Gauss
seminorm at its coordinate radii, which identifies the point.
The strict inequality persists on an ordinary neighborhood $U$.
At every point of $U$ at least two monomials realize the maximum,
since a unique largest term cannot cancel in a non-Archimedean
seminorm. Thus the coordinate tropicalization of $U$ lies in the
finite tropical hypersurface of $P$, of dimension at most one.
The pullback of a coordinate $(2,2)$ superform vanishes there:
its restrictions to the local tropical polyhedra have dimension
less than two. This is the dimension rule of the actual CLD
moment-form calculus. In particular $\theta^2$ has no support
off the closed coordinate skeleton.

On that skeleton the identity moment has saturated lattice and
finite-flat degree one. Its calibration is one by CLD14.5.2--3.
CLD5.10.1--2 identifies the integral with its calibrated tropical
integral; the raw Hessian-to-wedge factor is $2!$.
The Hessian of $|r|^2/2$ is the identity, so the tropical density
is $2\,dr_1dr_2$. Support on the skeleton and its Gauss-section
homeomorphism identify the full Radon measure, not only its
total mass. These identities glue through the quotient charts;
the fundamental square has area one. This proves the assertion.
\end{proof}

\section{Uniform mixed moment bounds from the bounded dual}
\begin{lemma}\label{mixed}
If $p,q,r\in\C$, all three are nonpositive and have energy at
least $-C_0$, then
\[
 \int(-p)\MA(q)\le3C_0,\qquad
 \int(-p)B(q,r)\le6C_0.
\]
Moreover, for nonpositive $a,w$ with the same energy bound,
\[
 Q_{\theta_w}(a,a)\le6C_0,\qquad Q_\theta(a,a)\le6C_0.
\]
\end{lemma}
\begin{proof}
All three terms in the definition of $-3E(q)$ are nonnegative,
so $\int(-q)\MA(q)\le-3E(q)$. The full bounded dual formula gives
\[
 E^*(\MA(q))=E(q)+\int(-q)\MA(q)\le-2E(q)\le2C_0.
\]
Its defining inequality for $p$ gives the first bound.
For $b=(q+r)/2$, concavity gives $E(b)\ge-C_0$ and $b\le0$.
Positive product expansion gives $B(q,r)\le2\MA(b)$, proving
the second bound. Finally Green gives
\[
 Q_{\theta_w}(a,a)=\int(-a)B(a,w)-\int(-a)B(0,w)
 \le\int(-a)B(a,w)\le6C_0.
\]
The $\theta$ assertion is the case $w=0$.
\end{proof}

\section{Fixed gradient measures have bounded mixed-measure domination}
For $T=\theta$ or $T=\theta_w$, define the finite positive
Radon gradient measure $\lambda_{w,T}$ by
\[
 \int\chi\,d\lambda_{w,T}=Q_{T,\chi}(w,w)
 \quad(\chi\text{ continuous}),\qquad
 \lambda_{w,T}=d'w\wedge d''w\wedge T
\]
in the completed-gradient meaning of the appendix. Positivity
and the bound by $\|\chi\|_\infty Q_T(w,w)$ give its actual
Radon interpretation by Riesz representation. This notation
does not presume smoothness of $w$.

\begin{lemma}\label{gradom}
For fixed $w\in\C$ and either such $T$, there are finitely many
bounded $v_j\in\C$ and positive constants $b_j$ with
\[
 \lambda_{w,T}\le\sum_jb_j\theta_{v_j}\wedge T.
\]
Consequently, for every $C\ge0$ there is a finite $H_T(C)$,
depending only on $w,T,C$, such that
\begin{equation}\label{H}
 \int(-a)\,d\lambda_{w,T}\le H_T(C)
 \quad(a\in\C,\ a\le0,\ \sup_Xa=0,\ E(a)\ge-C).
\end{equation}
\end{lemma}
\begin{proof}
Take finitely many injective tropical box charts with nested
inner boxes covering $X$ and outer boxes having compact closure.
On each choose a centered quadratic reference $\rho_0$, with
$\DD\rho_0=\theta$. By adding finite constants make
$\phi=\rho_0+w+B$ and $\psi_0=\rho_0+B_0$ at least one
on the entire outer closure. They are bounded continuous psh
there, at their actual sheaf stages. Their squares are psh by
convex nondecreasing composition on this range. The strong
gradient product rule gives
\[
 \DD(\phi^2T)=2\phi\DD(\phi T)+2\lambda_{\phi,T},
 \qquad \lambda_{\phi,T}\le\tfrac12\DD(\phi^2T).
\]
The same holds for $\psi_0$. Since $dw=d\phi-d\psi_0$,
the Hilbert inequality gives the local positive-measure bound
\[
 \lambda_{w,T}\le2\lambda_{\phi,T}+2\lambda_{\psi_0,T}
 \le\DD(\phi^2T)+\DD(\psi_0^2T).
\]
The rules follow by strong smooth approximation and bounded
continuous multiplication as in the appendix, so hold for every
continuous stage, not only smooth potentials.

We globalize each bounded local psh function $\psi$ appearing
on the right. Choose the inner box sufficiently small that the
maximum $b_0$ of $\rho_0$ there is smaller than its minimum
$b_1$ on an outer boundary collar. Choose $\epsilon>0$ so small
that $\epsilon(\max\psi-\min\psi)<b_1-b_0$ on the outer
closure, and choose $\delta$ between
$b_0-\epsilon\min\psi$ and $b_1-\epsilon\max\psi$.
Then the globally glued potential
\[
 v=\max(0,\delta+\epsilon\psi-\rho_0)
\]
on the box, extended by zero outside it, belongs to $\C$.
Indeed $\rho_0+v=\max(\rho_0,\delta+\epsilon\psi)$ is
actual continuous psh, and $v$ is identically zero on the outer
collar. On the inner box $\theta_v=\epsilon\DD\psi$.
Thus the preceding local bound is a positive combination of
actual bounded mixed products there. Summing over the finite
inner cover bounds the whole positive gradient measure by their
global sum, proving the first assertion. Normalizing each $v_j$
and $w$ by finite constants does not change their curvature or
gradient measures. Lemma\ref{mixed} gives\eqref{H}, for example
with $H_T(C)=6\sum_jb_j\max(C,-E(v_j),-E(t),0)$, where
$t=0$ or the fixed normalized $w$ according to $T$.
\end{proof}

\section{Uniform negative-part tails for continuous weights}
Normalize the fixed bounded test by $w\le0$, set
$C_w=\max(0,-E(w))$, and for $C\ge0$ put
\[
 K_w(C)=\sqrt{6\max(C,C_w)}
       \bigl(\sqrt{H_{\theta_w}(C)}+\sqrt{H_\theta(C)}\bigr).
\]
All these constants depend on the fixed test and the energy
bound; they do not depend on the varying potential or truncation.

\begin{theorem}\label{tailC}
For every $a\in\C$ with $\sup_Xa=0$ and $E(a)\ge-C$, and
every $T>1/2$,
\begin{equation}\label{tailineq}
 \int_X(-a-T)_+\MA(w)\le\frac{K_w(C)}{\sqrt T}.
\end{equation}
Consequently
\[
 \int_{\{a<-2T\}}(-a)\MA(w)\le\frac{2K_w(C)}{\sqrt T}.
\]
\end{theorem}
\begin{proof}
Put $d=(-a-T)_+=\max(a,-T)-a$, an actual continuous dpsh
function. By\eqref{refbound} it vanishes on $S$ when $T>1/2$.
Proposition\ref{refmeasure} gives $\int d\theta^2=0$.
Product transgression and Green therefore give
\begin{equation}\label{cross}
 \int d\MA(w)=-Q_{\theta_w}(d,w)-Q_\theta(d,w).
\end{equation}

Here is the localized estimate, including its nonsmooth
truncation justification. Choose smooth one-variable functions
$F_\epsilon$ uniformly within $\epsilon$ of
$q\mapsto(-q-T)_+$, with $|F_\epsilon'|\le1$ and
$F_\epsilon'=0$ for $q\ge-T+\epsilon$. Smooth composition
in the strong-gradient domain gives
$d'F_\epsilon(a)=F_\epsilon'(a)d'a$.
Cauchy--Schwarz with the bounded continuous multiplier
$F_\epsilon'(a)$ on the fixed gradient of $w$ gives
\[
 |Q_R(F_\epsilon(a),w)|\le
 \sqrt{Q_R(a,a)}
 \sqrt{\int|F_\epsilon'(a)|^2\,d\lambda_{w,R}}
 \le\sqrt{Q_R(a,a)}
 \sqrt{\lambda_{w,R}(\{a\le-T+\epsilon\})},
 \quad R\in\{\theta,\theta_w\}.
\]
Only the integral inequality compares this continuous multiplier
with an indicator; the gradient module is multiplied solely by
a bounded continuous function, as already established in the
appendix.

By\eqref{H}, the second squared factor is at most
$H_R(C)/(T-\epsilon)$ for $0<\epsilon<T$.
Lemma\ref{mixed} bounds the first squared factor by
$6\max(C,C_w)$. Since $\DD(wR)$ is signed Radon, Green
also gives
$Q_R(F_\epsilon(a),w)=-\int F_\epsilon(a)\DD(wR)$.
Uniform convergence therefore passes this bilinear pairing to
$Q_R(d,w)$; no unsupported strong convergence of discontinuous
indicator derivatives is required. Let $\epsilon$ tend to zero
and use\eqref{cross}. Its left side is nonnegative, and the two
absolute bounds sum to\eqref{tailineq}. On $\{a<-2T\}$,
$-a\le2(-a-T)_+$, giving the final assertion.
\end{proof}

\section{The bound passes to actual decreasing nets and poles}
We record the actual topological integral fact used in this
passage. If $f_i\ge0$ is an increasing net of continuous
functions on compact Hausdorff $X$, and $\mu$ is finite Radon,
then
\begin{equation}\label{radonnet}
 \int\sup_i f_i\,d\mu=\sup_i\int f_i\,d\mu.
\end{equation}
To see this, $f=\sup_i f_i$ is lsc. Radon inner regularity on
each open $\{f>k\delta\}$ and compact-Hausdorff Urysohn
functions produce continuous $0\le h\le f$ whose integrals
approximate its finite layer-cake sums. Explicitly take compact
$K_k$ in those open sets, choose continuous $h_k\in[0,1]$,
one on $K_k$ and supported in $\{f>k\delta\}$, and use
$h=\delta\sum_{k=1}^N h_k\le f$. Thus $\int f$ is the
supremum of integrals of continuous such $h$.
For each $h$ and $\epsilon>0$, the open sets
$\{f_i>h-\epsilon\}$ cover $X$, so one index works everywhere
by a finite subcover and directedness. This proves\eqref{radonnet},
also for infinite integrals. It does not assert monotone
convergence for arbitrary measurable nets.

\begin{theorem}\label{tailN}
If $u\in\GN$, $\sup_Xu=0$ and $E(u)\ge-C$, then
\[
 -\tfrac12\le u|_S\le0,\qquad
 \int(-u-T)_+\MA(w)\le\frac{K_w(C+1)}{\sqrt T}
 \quad(T>1/2).
\]
In particular the normalized net-generated energy sublevel is
uniformly integrable against $\MA(w)$, for every fixed $w\in\C$.
Every such $u$ is finite $\MA(w)$-almost everywhere.
\end{theorem}
\begin{proof}
Fix an actual continuous decreasing presentation $a_i\downarrow u$.
Its energy infimum is canonical: compact strict domination
compares any strict continuous majorant with a presentation term;
adding an arbitrary positive constant to a term gives the reverse
comparison. Thus $E(u)=\inf_iE(a_i)$, as in406, and
$E(a_i)\ge-C$. The open sets $\{a_i<\epsilon\}$ cover $X$
because $u\le0<\epsilon$; one index works globally by
compactness. Hence $s_i=\sup_Xa_i\downarrow0$ and $s_i\ge0$.
Apply Theorem\ref{gauss} to $b_i=a_i-s_i$ to get
$a_i|_S\ge s_i-1/2$. Passing to the limit proves the reference
bound, and the Lipschitz bounds also pass pointwise.

On a tail $0\le s_i\le1/2$, the normalized $b_i\in\C$ has
$E(b_i)=E(a_i)-2s_i\ge-C-1$.
Since $a_i\ge b_i$, Theorem\ref{tailC} gives
\[
 \int(-a_i-T)_+\MA(w)\le\int(-b_i-T)_+\MA(w)
 \le K_w(C+1)/\sqrt T.
\]
The continuous nonnegative functions $(-a_i-T)_+$ increase to
$(-u-T)_+$ on this tail. Use\eqref{radonnet} for the whole
finite positive test measure to prove the asserted inequality,
including at poles. The final tail bound follows as before by
$-u\le2(-u-T)_+$ on $\{u<-2T\}$.
Since the right side tends to zero, it gives uniform
integrability. The finite truncation integral also rules out
positive measure of any pole set.
\end{proof}

\section{Exact scope and remaining compactness gates}
The new proofs above supply a fixed-bounded-test uniform tail
estimate, stronger than an $L^1$ bound. Finite sums of bounded
mixed measures inherit it because $B(v,w)\le2\MA((v+w)/2)$.
It is not a uniform tail theorem for every finite-full-dual
measure with no additional domination. Nor does it give an actual
positive presentation for an arbitrary weak or Shilov-coordinate
limit. Upper-envelope regularity, arbitrary-limit representability,
actual energy-sublevel compactness and arbitrary-support
Monge--Ampere existence remain OPEN. None of these is assumed
in the truncation proof. A countable dense set or a full-support
finite Radon measure on this nonmetrizable space is not presumed.

The complete immediate GRv2 Proposition4.2 statement/proof,
its annulus/affinoid maximum Theorem6.4, and the unnumbered
Section1.2 smooth recognition argument retain their exact
earlier reading qualifications. New immediate CLDv2 readings
include the whole5.10.1 statement/proof and5.10.2 reconstruction
at LF6936--6995, and the whole14.5.2--3 calibration argument at
LF20953--21008. The dimension rule for moment pullback forms,
explicit Gauss norms and annulus Shilov boundaries are named
conventional CLD/Berkovich foundations at these torus hypotheses.
No whole recursive audit of every referenced tropical-volume,
analytic-base-extension or curve-potential theorem is invented.

The full current402 mathematical body is reproduced below as
the actual bounded-product/strong-gradient/energy foundation.
Its exact identity is retention, not another independent proof
reading or an incoming correctness grant. The essential new
proofs here are complete; all mathematical claims are
UNREVIEWED. Compilation and identities do not certify them.
Full primary bodies/literal units/pixels/book contexts remain
LOCALONLY. Counts16/42 solved,26OPEN,97accepted scopedpapers
and original closure are unchanged. The protected P16.4 editor
remains OPEN UNCHANGED, and this work is not sent as hot
feedback to active frozen campaigns. Native/generated programs
remain DATA and are never executed or imported.

\appendix
\section{Complete current402 bounded-product and energy foundation}
The following complete mathematical body is preserved literally.
Its section structure and labels are unchanged.
\section{The actual analytic object and proof foundations}
Put $K=\mathbb C((s))$, $|s|=e^{-1/8}$, $t=s^8$,
$Y=(\mathbb G_m^2)^{\mathrm{an}}_K$, and
\[
 X=Y/\langle(t,2),(3,t)\rangle,\qquad
 \tau=(-\log|z_1|,-\log|z_2|).
\]
The translations on $\tau$ form $\mathbb Z^2$. Ordinary open tropical
boxes of width less than one give injective quotient charts. We use
the previously reconstructed analytic quotient: compact, separated,
good, boundaryless, connected, smooth, and pure of dimension two.
The quotient construction and its nonalgebraization proof are not
reproved here and are not consequences of the present point test.
On a chart let $\rho=|\tau|^2/2$ and $\theta=\DD\rho$.
Changes of reference are affine unit logarithms. With the identity
torus's saturated lattice and degree one,
\[
 M=\int_X\theta^2=2.
\]
This is the raw CLD convention; no probability normalization is made.

Let $\mathscr P$ be the minimal continuous CLD psh sheaf containing
the ordinary smooth psh functions and closed under local uniform
limits. All stages of this sheaf are retained. Set
\[
 \C=\{a\in C(X,\RR):\rho+a\in\mathscr P\text{ locally}\},\quad
 \theta_a=\theta+\DD a,\quad B(a,b)=\theta_a\wedge\theta_b.
\]
We use the actual recursive continuous products. Their positivity,
symmetry, closedness, locality, local uniform continuity, and
invariance under smooth $\DD$-null reference changes are precisely
CLDv2 8.5.1, 8.5.3, 8.5.5--6, 8.7.1--4 and 8.8.4--5, at the stated
continuous psh/dpsh and positive closed-current hypotheses. Positive
top currents are Radon measures by 7.2.8 on these locally
differentially separated charts. Smooth cutoffs, partitions, scalar
approximation and boundaryless Stokes are the conventional analytic
foundations. Arbitrary continuous-multiplier closedness in 8.6.7 is
not a premise. No first-stage identity for $\mathscr P$ is assumed.

Constants, convex combinations and finite maxima preserve $\C$.
For a smooth slot, convex regularized maxima are smooth psh.
Uniform continuity and sheaf minimality extend the operation one
slot at a time to all stages. This also proves ordinary convex
composition with finite continuous psh slots by the same argument.

\section{Bounded Green calculus and the full bounded dual}
Here we give the extra energy argument rather than infer it from
a positivity label. For a positive closed $(1,1)$ current $T$,
complete compact smooth $(1,0)$ forms modulo the null space of
\[
 (a,b)\longmapsto\int a\wedge Jb\wedge T.
\]
Use a nonnegative compact smooth cutoff $\chi$ for the corresponding
local seminorm $Q_{T,\chi}$. Let $\HH_T$ consist of continuous
functions locally uniformly approximable by smooth functions whose
gradients converge strongly in these seminorms. The gradient is
unique: for a uniformly zero sequence $f_j$ with a strong gradient
limit and a compact smooth form $b$, closedness gives
\[
 \int d'f_j\wedge Jb\wedge T=-\int f_jd'(Jb)\wedge T\longrightarrow0.
\]
The right-hand coefficient measure is finite on the fixed compact
support; smooth Hilbert tests are dense. The limit gradient is zero.
Partitions patch local approximants by the Leibniz rule.
Bounded continuous coefficients act by scalar approximation and the
bound $\|fa\|\le\|f\|_\infty\|a\|$. Products and cutoffs have their
usual gradient rules.

If $f\in\HH_T$ and $L_Tf=\DD(fT)$ is signed Radon, smooth
approximation first in $f$, then uniformly and strongly in a
compactly supported $g$, yields
\[
 \int gL_Tf=-Q_T(f,g),\qquad
 Q_{T,\chi}(w,w)=\frac12\int w^2\DD\chi\wedge T-\int\chi wL_Tw.
\]
The sheaf of $f\in\HH_T$ with $L_Tf\ge0$ contains smooth psh
functions and is uniformly closed. Indeed take nested cutoffs
$\chi,\eta$, $\eta=1$ near $\operatorname{supp}\chi$,
a common bound $|f_j|\le B$ on the whole support of $\eta$, and
$\epsilon=\|f_j-f_k\|$ there. Positivity and the preceding identity give
\[
 0\le Q_{T,\chi}(f_j-f_k,f_j-f_k)
 \le\frac{\epsilon^2}{2}\|\DD\chi\wedge T\|_{\mathrm{TV}}
 +2B\epsilon\|\chi\|_\infty\|\DD\eta\wedge T\|_{\mathrm{TV}}.
\]
Choose smooth approximants to $f_j$ with uniform and smaller-patch
gradient errors below $1/j$. The estimate proves strong convergence;
distributional positivity survives. Minimality places every stage
of $\mathscr P$ in this sheaf. Its vector span includes continuous
dpsh functions. Finite patching on compact $X$ proves
\[
 \int f\DD(gT)=\int g\DD(fT)=-Q_T(f,g),\qquad Q_T(f,f)\ge0.
\]
In particular all the following computations hold at every
continuous stage. The exact terms in the bounded product expansion
have zero total integral, so every $B(a,b)$ has mass $M=2$.

Define
\[
 E(a)=\frac13\int_X a[\theta^2+\theta_a\wedge\theta+\theta_a^2].
\]
For $a_t=(1-t)a+tb$, the cubic expansion and Green give
\[
 \frac{d}{dt}E(a_t)=\int(b-a)\MA(a_t),\qquad
 \frac{d^2}{dt^2}E(a_t)=-2Q_{\theta_{a_t}}(b-a,b-a)\le0.
\]
Thus $E$ is concave and monotone, $E(0)=0$, and
$E(a+c)=E(a)+Mc$. Its tangent inequality is
\[
 E(a)\le E(w)+\int(a-w)\MA(w)\qquad(a,w\in\C).
\]
For a positive Radon measure $\mu$ define the \emph{full} bounded dual
\[
 E^*(\mu)=\sup_{a\in\C}\left(E(a)-\int a\,d\mu\right).
\]
Constants force $\mu(X)=M$ whenever this quantity is finite.
The tangent inequality, with equality at $a=w$, proves the exact
formula
\begin{equation}\label{dual}
 E^*(\MA(w))=E(w)-\int w\MA(w)<\infty
 \quad\text{for every }w\in\C.
\end{equation}
It is nonnegative, since $a=0$ is one of the competitors.

\section{The actual generated class and a moment bound}
Let $\G$ consist of actual pointwise limits
$a_j\downarrow u$ on $X$, where $a_j\in\C$ and
$\inf_jE(a_j)>-\infty$. These limits are USC.
They are not abstract functionals or equivalence classes modulo
null sets. Define $E(u)=\lim_jE(a_j)$.
This is independent of every continuous decreasing presentation.
To see this, if $b_k\downarrow u$ and $a$ is continuous with $u<a$
everywhere, the increasing ordinary open sets $\{b_k<a\}$ cover
compact $X$. One index therefore works everywhere. Apply this to
$a=a_j+\epsilon$, use monotonicity and the constant rule, and
then reverse the roles of the presentations. The limits agree.
In particular $E(u)>-\infty$ for every presentation of a member
of $\G$. The identically $-\infty$ limit is excluded: compact
strict domination would eventually put a presentation below every
negative constant, contradicting its energy lower bound.

\begin{lemma}\label{moment}
If $u\in\G$, $u\le0$, and $E^*(\mu)<\infty$ for a positive
Radon measure $\mu$, then
\[
 \int(-u)\,d\mu\le E^*(\mu)-E(u)<\infty.
\]
\end{lemma}
\begin{proof}
Every term of a presentation satisfies
$\int(-a_j)\,d\mu\le E^*(\mu)-E(a_j)$.
The continuous first term bounds all terms above. Apply monotone
convergence to that common bound minus $a_j$ and pass to the limit.
This proves the claim even at poles and does not presume
integrability before the bound. The same argument proves
$L^1(\mu)$ convergence of the whole presentation.
\end{proof}

\section{Laurent Abhyankar moments without changing the field}
Write $d_K(x)=\operatorname{trdeg}_{\widetilde K}
\widetilde{\mathscr H(x)}$, with graded residue fields.
Let $\mathscr A$ be the set of points of $X$ with $d_K(x)=2$.
A lift of such a point in $Y$ will also be denoted $x$.

\begin{lemma}\label{moments}
There are two nonzero Laurent polynomials $P_1,P_2$ over $K$
whose initial forms at $x$ are algebraically independent over
$\widetilde K$. Put $r_i=|P_i(x)|>0$. Then
\[
 (P_1,P_2)(x)=\eta_{(r_1,r_2)}
\]
in $(\mathbb A_K^2)^{\mathrm{an}}$. The map is finite and flat at
$x$, and admits a local finite flat representative with a positive
constant degree $D$ and with $x$ the only point above this Gauss point.
\end{lemma}
\begin{proof}
The support of $x$ on the Laurent polynomial ring is zero.
Otherwise it lies on a proper algebraic subvariety of dimension
at most one, contradicting $d_K(x)=2$ and the Abhyankar inequality.
Thus $\mathscr H(x)$ is the completion of the valued field
$K(z_1,z_2)$. Completion does not add graded residue elements:
approximate a nonzero element with error strictly less than its
norm. Its initial form agrees with that of the approximant.
Consequently the initial forms of all nonzero Laurent polynomials
generate the graded residue \emph{field}, allowing ratios.
A maximal algebraically independent subfamily of these homogeneous
generators has cardinal two. Choosing its members gives $P_1,P_2$.

For every polynomial $\sum_I c_IP^I$, the initial terms of maximal
norm cannot cancel by graded algebraic independence. Its norm is
$\max_I|c_I|r^I$. This is exactly the Gauss seminorm, for arbitrary
real $r_i>0$; no scaling of $r_i$ into $|K^\times|$ is required.
This is also the Abhyankar-basis criterion of CLD 13.8.5.

The torus is boundaryless, hence $x$ is in the relative interior
of the map to $\mathbb A^2$. The target is reduced and the Gauss
point has graded dimension two. CLD 11.6.5(c) makes the map finite
and flat at $x$. Its 11.6.3 local representative can be chosen
with a single point of that fibre. The rank is a positive integer
$D$. After shrinking the target inside its ordinary interior,
the source is an ordinary open neighborhood of $x$, and the finite
flat degree remains $D$ there.
\end{proof}

\section{The actual atomic mass and its global test}
We spell out the local calibrated number.
On a real two-plane set $g(t)=\max(0,t_1,t_2)$ and
$g_\epsilon=\epsilon\log(1+e^{t_1/\epsilon}+e^{t_2/\epsilon})$.
For $g_1$, write $p_i=e^{t_i}/(1+e^{t_1}+e^{t_2})$ and
$p_0=1-p_1-p_2$. The Hessian determinant is $p_0p_1p_2$.
The gradient maps the plane diffeomorphically to
$\{p_1,p_2>0:p_1+p_2<1\}$ with that Jacobian.
The raw superform rule is $2!\det(\operatorname{Hess}g_1)\,dt_1dt_2$.
Its integral is $2!\,(1/2)=1$.
Dilation and dominated convergence concentrate this finite measure
at the origin as $\epsilon\downarrow0$, so
\[
 (\DD g)^2=\delta_0.
\]
This is the explicit case of CLD 1.7.1 and 1.7.6--8, including
their factor $2!$.

Near $x$, both $P_i$ are invertible. Apply the pleasant moment
$(P_1,P_2)$ there. Finite flatness and the single Gauss fibre imply
full tropical dimension and that the calibrated tropical cone at
$x$ is the whole real plane with multiplicity $D$.
Indeed the map is open, so the Gauss section of a neighborhood of
the target gives a neighborhood in the tropical plane. The explicit
calibration formula CLD 14.5.2--3 makes its image calibration
$D$ times the standard one. Locality removes other fibre points,
since there are none in this representative. More explicitly,
smooth log-sum-exp approximants have top measures strongly
supported on the inverse image of the target Gauss section.
The calibrated integration formula pushes them to $D$ times the
corresponding real-plane measures near the target tropical value.
Those measures concentrate there with mass one. Finite properness
and the single Gauss fibre imply that every inverse Gauss point
over a sufficiently small tropical neighborhood lies in any
prescribed ordinary neighborhood of $x$: otherwise compactness
would give another point of that fibre. Continuous tests therefore
converge to $D$ times evaluation at $x$. Local uniform continuity
of the actual product identifies the limit with the log-maximum
product. This reconstructs the atomic calculation 8.9.2--4 and yields
\begin{equation}\label{localatom}
 \left(\DD\max(0,\log(|P_1|/r_1),\log(|P_2|/r_2))\right)^2
 =D\delta_x
 \quad\text{on some ordinary neighborhood of }x.
\end{equation}
Formula \eqref{localatom} uses the normalized
canonical calibration over the unchanged field $K$ and the actual
finite flat degree. No arbitrary-current pullback is invoked.
CLD 18.6.2 gives the equivalent Gauss-fibre description, but its
relative-intersection proof is not needed for this argument.

\begin{theorem}\label{atom}
For every $x\in\mathscr A$ there is $w_x\in\C$ and
$\alpha_x>0$ such that the actual positive Radon measure
$\MA(w_x)$ satisfies
\[
 \MA(w_x)(\{x\})=\alpha_x.
\]
One can take $\alpha_x=\varepsilon^2D$ in the construction below.
\end{theorem}
\begin{proof}
Choose an injective lift chart around the tropical class
$a=\tau(x)$ and a small Euclidean tropical ball of radius $R$ there.
Its inverse image has compact closure inside a polyannulus.
Set $\rho_a=|\tau-a|^2/2$; its difference from $\rho$ is
smooth affine and $\DD$-null.
The function
\[
 \psi=\max(0,\log(|P_1|/r_1),\log(|P_2|/r_2))
\]
is finite, continuous, nonnegative and psh on this whole chart,
including zeros of the polynomials. Here is a direct check at all
continuous stages. Choose a smooth convex $\chi$ equal to zero for
$v\le0$, with $0\le\chi'\le1$, and equal to $v-c$ for $v\ge1$,
where $0<c<1$. For $N>3$ the finite smooth psh function
\[
 p_{i,N}=\chi(\log(|P_i|/r_i)+N)-N+c
\]
is locally constant at zeros. Where it can be positive it equals
$\log(|P_i|/r_i)$, so $\psi=\max(0,p_{1,N},p_{2,N})$ exactly.
Smooth convex regularized maxima place it in $\mathscr P$.

Take $0<\delta<R^2/4$ and choose $\varepsilon>0$ so small that
$\delta+\varepsilon\sup\psi<R^2/2$ on the compact closure.
Then
\[
 w_x=\max(\rho_a,\delta+\varepsilon\psi)-\rho_a
\]
vanishes on a tropical boundary collar. Extend it by zero to $X$.
It belongs to $\C$: regularize the inner finite maximum, and then
use a smooth convex regularized outer maximum which agrees exactly
with the larger slot when the gap exceeds its regularization
width. The fixed collar gap ensures that each sufficiently small
regularization extends by zero smoothly and $\theta$-psh.
Their uniform limit is the displayed function.

On the fixed inner neighborhood $\rho_a<\delta/2$ we have
$\theta_{w_x}=\varepsilon\DD\psi$. Shrink further to the
neighborhood in \eqref{localatom}. Locality of the actual recursive
product gives
$\MA(w_x)=\varepsilon^2D\delta_x$ there.
Outside it the global product remains positive and has raw mass
$M=2$. Thus its atom at $x$ is precisely the positive number
$\alpha_x=\varepsilon^2D$.
\end{proof}

\section{Finiteness, sublevel bounds, and the precise compactness limit}
\begin{corollary}\label{finite}
For every $u\in\G$ and every $x\in\mathscr A$, $u(x)$ is finite.
More precisely, if $\sup_Xu=0$ and $E(u)\ge-C$, then
\[
 -\frac{E^*(\MA(w_x))+C}{\alpha_x}\le u(x)\le0 .
\]
\end{corollary}
\begin{proof}
Use \eqref{dual} and Lemma \ref{moment} for the \emph{whole}
positive measure $\mu_x=\MA(w_x)$.
Since $-u\ge0$ and $\mu_x$ contains the positive atom
$\alpha_x\delta_x$,
\[
 \alpha_x(-u(x))\le\int(-u)\,d\mu_x
 \le E^*(\mu_x)-E(u)\le E^*(\mu_x)+C.
\]
This excludes a pole at $x$ and proves the bound.
An arbitrary $u$ can be shifted by its finite supremum.
A member of $\G$ has a finite upper bound and is not identically
$-\infty$, so that shift is legitimate.
\end{proof}

Every Shilov point of a pure-dimensional two-dimensional affinoid
domain in a local chart has $d_K=2$. At the exact general-radius
hypotheses, GRv2 Lemma6.3(3)--(4) expresses these Shilov points as
the strata with two graded algebraically independent generators;
the Abhyankar inequality gives the reverse bound.
The entire lemma and its immediate proof were read, including its
Temkin graded-reduction and Ducros fibre-dimension foundations.
These deeper theorems are named published foundations, not freshly
recursively reproved here. Hence Corollary \ref{finite} applies to
every such Shilov point, not only to the reference torus skeleton.
On a fixed finite Shilov boundary the lower bound can be chosen
uniformly by taking the maximum of the finitely many displayed
constants.

For clarity, actual local weights are classically psh at all
stages. The published smooth CLD-to-classical bridge is the
unnumbered statement and argument in GRv2 Section1.2.
Its Proposition4.2 gives sheafiness and decreasing-net
closure. A uniform limit of finite continuous classical psh
functions is classical psh: select errors at most $2^{-j}$ and add
$3\cdot2^{-j}$, obtaining a decreasing sequence with the same limit.
Sheaf minimality then recognizes all of $\mathscr P$, and a
continuous decreasing presentation recognizes $\rho+u$ at poles.
GRv2 Theorem6.4 therefore gives the affinoid Shilov maximum
principle for these actual USC weights.
For all arbitrarily small affinoid neighborhoods $A$ in one chart,
\[
 (\rho+u)(y)=\inf_{A\ni y\ \mathrm{neighborhood}}
                \max_{\xi\in\Gamma_A}(\rho+u)(\xi).
\]
The maximum principle gives one inequality. USC gives the other
by every finite upper bound at a pole and by
$(\rho+u)(y)+\eta$ at a finite value. Thus values on the union
$\mathscr S$ of \emph{all} these Shilov boundaries determine the
actual weight, including all poles.

Fix a normalized energy sublevel
$\G_C=\{u\in\G:\sup u=0,\ E(u)\ge-C\}$.
Its evaluation map on $\mathscr S$ is injective and takes values in
\[
 \prod_{\xi\in\mathscr S}[-L_\xi(C),0],\qquad
 L_\xi(C)=\frac{E^*(\MA(w_\xi))+C}{\alpha_\xi}.
\]
Tychonoff compactness implies that every net of such evaluations
has a coordinatewise convergent subnet with finite coordinates.
This is compactness in the product space of boundary values.
It does \emph{not} show that the limit coordinates are represented
by an actual member of $\G$, that the normalized image is closed,
or that energy is semicontinuous there. Those are still mathematical
gates for actual energy-sublevel compactness.
Likewise this point bound alone does not prove uniqueness from
equality of generated MA measures.

\begin{remark}
We have not claimed that $\delta_x$ itself has finite full bounded
dual merely because it is part of $\MA(w_x)$.
The whole test measure has the supporting-energy formula
\eqref{dual}; its positivity is exactly what gives the atomic
point inequality. The distinction retains the raw mass-two dual
and avoids an invalid componentwise assertion.
\end{remark}

\section{Source scope and original residual}
The exact CLDv2 full source is bound locally by SHA256
\texttt{64846cb00926267c6238d1afac61b5eaed0a0e3a6349725738f9b1bf163de085}.
Applied immediate readings include 1.7.1 and1.7.6--8
(LF1773--1830 and1931--2135), the complete local finite-flat block
11.6.1--5 (LF15846--15912), the graded Abhyankar criterion13.8.1--5
(LF19772--19820), the full calibration formula/proof14.5.2--3
(LF20953--21008), and the entire atomic calculation8.9.1--4
(LF10595--10739). The LF windows actually retained may include
adjacent context; essential statements and whole immediate proofs
were read in bounded outputs. The full relative atom formula18.6.2
(LF27607--27671) was read as corroboration, not a premise of the
finite-flat/calibration proof.
No algebraically closed standing hypothesis from14.6 is imported.
The real-plane coefficient one was computed above rather than
inferred from an OCR inequality.

The continuous products and bounded Green foundations retain the
exact earlier 8.5--8.8 source applications and all their hypotheses;
the complete fresh point construction and energy argument are
given here. GR is the full17September2025 v2 source, bound by
\texttt{3a221b86116604c950db07e4c60ab184701f0624deebabd92502e5c81871436b}.
Lemma6.3 and its whole immediate proof were read at LF1673--1780.
Its deeper Temkin/Ducros reduction proofs, basic Berkovich
relative-interior/dimension and completed-residue facts, graded
field transcendence, and the analytic quotient construction are
disclosed conventional foundations. There is no claim of an
entire new recursive source audit. All full primary bodies,
literal units, pixels and book contexts remain LOCALONLY.

Original P16.7 still requires actual separated energy-sublevel
compactness/representability, envelopes or a proved variational
replacement, arbitrary-support equation existence at the stated
full bounded-dual criterion, and full generated pointwise uniqueness.
This paper proves the previously missing actual point-finiteness
and coordinate lower bounds. It neither closes the original nor
changes any accepted count. The protected P16.4 editor remains open
and unchanged; new work is not sent as hot feedback to active
native campaigns. All claims remain UNREVIEWED.

\end{document}
